\chapter{Context Engineering}
\label{ch:context}

\begin{quote}\itshape
The context window is a budget you are spending whether or not you are managing
it.
\end{quote}

\noindent
The source material this book grew from marks this as its largest missing
chapter, and it is right to. Everything in Parts~II and~III assumes the model can
see what it needs to see, and in a real agent that assumption fails quietly:
tool outputs accumulate, history grows, quality degrades before any limit is hit,
and the bill rises the whole time. Context engineering is the discipline of
treating the window as the scarce, managed resource it is.

\chapterstub{
\textbf{Argument.} Tool results, not conversation history, are the dominant
consumer of context in a working agent, and they are the thing nobody trims. A
single unfiltered SQL result or web page can cost more than the entire preceding
conversation.

\textbf{Sections.} (12.1) The window as a managed resource: what occupies it, in
what proportion, measured on a real run rather than assumed. (12.2) Compaction
and summarisation: the summarisation middleware, what it keeps, and what it
silently loses. (12.3) Trimming tool results --- the main event: truncation
strategies, projecting only the fields the model needs, and returning a handle
rather than a payload. (12.4) Editing context in flight with the context-editing
middleware. (12.5) Prompt caching: where cache points go, why message ordering
determines whether they hit, and the economics with real numbers. (12.6) Quality
degradation over long contexts --- what the evidence actually supports, and the
practical rule that follows. (12.7) Tool selection under scale, the selector
middleware, and MCP namespacing. (12.8) Injecting memory and user profile without
paying for it on every turn.

\textbf{Listings.} A context accounting harness that prints the window
composition per turn; a tool-result projector; summarisation configured and its
loss demonstrated; a prompt-cache-friendly message layout against a
cache-hostile one, with the cost difference.

\textbf{Diagrams.} \code{ctx-window-budget} --- window composition over ten turns
as a stacked bar. \code{ctx-compaction} --- before and after summarisation.
\code{ctx-cache-points} --- why reordering messages destroys a cache hit.

\textbf{Measurement.} Cost per conversation across four configurations: naive,
trimmed tool results, summarised history, and both --- plus cache hit rate.
Ten conversations of ten turns each. This one is cheap and the numbers are
striking; it should be run early.

\textbf{Mini-project.} Stage 09: tool-result projection, summarisation, cache
points, and a cost report per conversation.

\textbf{Source topics covered.} Part~V 5.1--5.8, all marked as gaps in the source.
}
